\chapter{Définitions des Variables de Derrida et Gardner (1984)}\label{annexe:variables}

Cette annexe recense les variables et notations introduites dans les travaux de Derrida et Gardner~\cite{Derrida1984} concernant le développement en faible désordre de l'exposant de Lyapunov.

\begin{itemize}
    \item $\lambda$ : Paramètre utilisé pour le développement en faible désordre ($\lambda \rightarrow 0$).
    \item $E$ : Paramètre d'énergie fixe.
    \item $\gamma(E)$ ou $\gamma$ : Exposant de Lyapunov associé à un produit donné de matrices aléatoires.
    \item $\psi_n$ : Fonction d'onde évaluée au site $n$.
    \item $V_n$ : Potentiel aléatoire au site $n$.
    \item $\rho(V)$ : Distribution de probabilité donnée du potentiel aléatoire $V$.
    \item $U_n$ : Vecteur à deux composantes contenant les fonctions d'onde $\psi_{n+1}$ et $\psi_n$.
    \item $N$ : Limite supérieure dans le produit de matrices aléatoires ou nombre total de termes.
    \item $\epsilon$ : Paramètre nul pour les systèmes dynamiques intégrables et non nul pour les propriétés de mélange. Représente également une partie imaginaire infinitésimale ajoutée aux valeurs réelles de l'énergie.
    \item $R_n$ : Rapport des fonctions d'onde sur des sites adjacents, défini par $\psi_n / \psi_{n-1}$, mesurant la direction du vecteur $U_n$.
    \item $q$ : Paramètre réel utilisé pour définir une énergie $E = 2 \cos(qa)$.
    \item $A, B_n, C_n, D_n$ : Variables du développement en faible désordre de $R_n$ indépendantes de $\lambda$.
    \item $x$ : Paramètre mis à l'échelle définissant l'écart par rapport à un niveau d'énergie, tel que $E-2 = \lambda^{4/3}x$ près du bord de bande ou $E' = 2 \cos(\pi\alpha) + \lambda^2 x$.
    \item $P(R, E, \lambda)$ : Distribution de probabilité stationnaire pour la variable $R_n$.
    \item $\tilde{\rho}(E)$ : Densité d'états à l'énergie $E$.
    \item $t$ : Variable mise à l'échelle pour $R$, définie par la relation $R = 1 + \lambda^{2/3}t$.
    \item $H(t, x, \lambda)$ : Fonction de distribution de probabilité mise à l'échelle dérivée de $P$ via la transformation $H(t, x, \lambda) = \lambda^{2/3}P(1+\lambda^{2/3}t, 2+\lambda^{4/3}x, \lambda)$.
    \item $H_0, H_1, H_2$ : Fonctions successives dans le développement en série de $H(t, x, \lambda)$ par rapport à $\lambda$.
    \item $C, C_1$ : Constantes d'intégration arbitraires.
    \item $X$ : Paramètre mis à l'échelle défini comme $x / \langle V^2 \rangle^{2/3}$.
    \item $\alpha$ : Paramètre rationnel utilisé pour définir des points d'énergie spécifiques via $E = 2 \cos(\pi\alpha)$.
    \item $\varphi$ : Variable angulaire introduite par le changement de variables $R = \sin(\varphi+\pi\alpha)/\sin\varphi$.
    \item $G(\varphi)$ : Distribution de probabilité transformée définie par $P(R, E, \lambda) \frac{dR}{d\varphi}$.
    \item $G_0, G_1, G_2$ : Fonctions successives dans le développement en $\lambda$ de $G(\varphi)$.
    \item $E'$ : Valeur d'énergie perturbée au voisinage de $2 \cos(\pi\alpha)$ définie par $E' = E + \lambda^2 x$.
    \item $r, s$ : Entiers utilisés pour définir le nombre rationnel $\alpha = r/s$.
    \item $W(\varphi)$ : Fonction périodique de période $\pi/3$ apparaissant lors de la résolution des équations au voisinage de $E=1$.
\end{itemize}
